%%%%%%%%%%%%%%%%%%%%%%%%%%%%
%This is the IEEE Consumer Electronics Magazine's LaTeX template. Authors are requested to follow the instructions in the NOTEs.
%%%%%%%%%%%%%%%%%%%%%%%%%%%%

\documentclass{IEEEmce}

\usepackage[colorlinks,urlcolor=blue,linkcolor=blue,citecolor=blue]{hyperref}

\usepackage{hyperref}
\hypersetup{ 
 colorlinks=true, 
 linkcolor=red, 
 filecolor=blue, 
 citecolor = blue, 
 urlcolor=blue, 
 } 
 
\usepackage{upmath}
\usepackage{amsmath}

\jvol{1}
\jnum{1}
\paper{1}
\jmonth{Nov/Dec}
\publisheddate{25 11 2025}
\currentdate{01 12 2025}
\jname{IEEE Consumer Electronics Magazine}
\pubyear{2025}
\doiinfo{MCE.2025.000001}

\newtheorem{theorem}{Theorem}
\newtheorem{lemma}{Lemma}

\setcounter{secnumdepth}{0}

\begin{document}

%%%%
% NOTE: Add the Running Head Title below!
%%%%
\sptitle{Predicting Grades from Noisy Data} 

%%%%
% NOTE: The article title must be at most four lines.
%%%%

\title{Challenges of Grade Prediction: A Machine Learning Study on Noisy Synthetic Educational Data}

%%%%
% NOTE: Authors from the same institution in the following sequence must be listed on a single line! Do NOT indicate the Corresponding Author behind the authors' name. Any acknowledgment must go after the conclusion section.
%%%%

\author{Barna Berekméri}

%%%%
% NOTE: Indicate only the university/institution without departments/subdivision, without address, without country, without IEEE membership, or without email. Do not use abbreviations.
%%%%

\affil{Eötvös Loránd University}


%%%%
% NOTE: Add the Running Head and Article Titles below.
%%%%

\markboth{Predicting Grades from Noisy Data}{Challenges of Grade Prediction: A Machine Learning Study on Noisy Synthetic Educational Data}

%%%%
% NOTE: Do not break the line after \begin{abstract} command.
%%%%

\begin{abstract} This study investigates student performance prediction using a synthetic educational dataset containing over 1,000 records and multiple academic and behavioral features. A Random Forest regression model is then trained to estimate students’ final grades based on variables such as attendance and previous academic performance. Although the dataset’s synthetic nature results in very weak correlations between input features and the target variable, the analysis illustrates the challenges of predictive modeling in noisy environments and highlights the importance of feature relevance in machine learning workflows. The study also demonstrates how visualization and feature importance techniques support understanding model behavior and dataset limitations.\end{abstract}

\maketitle

\enlargethispage{10pt}

%%%%
% NOTE: Do not add an 'INTRODUCTION' section title. Instead, start the beginning of the first paragraph in \chapterinitial{} command; see below.
%%%%


%%%%
% NOTE: The main SECTION titles must be CAPITALIZED! Do not number the section/subsection titles.
%%%%

\section{Introduction} 

Predicting student performance has become an increasingly important topic in modern educational research, as accurate forecasts can support early intervention, personalized instruction, and resource allocation. The rise of machine learning has enabled new analytical approaches for modeling academic outcomes by learning patterns from historical data. Regression models, combined with feature analysis are widely used tools for understanding how various behavioral and academic factors contribute to students’ final results.

The aim of this study is to construct a complete data-processing and predictive modeling pipeline using a publicly available synthetic student performance dataset [forrás]. The workflow includes data cleaning, outlier correction, missing value imputation, categorical feature encoding, and the preparation of training and test sets. A Random Forest regression model is then trained to estimate students’ final grades, followed by a detailed evaluation of model accuracy through metrics and visualizations such as distribution plots, predicted-versus-actual comparisons, and feature importance analysis.

The dataset contains over 1000 generated student records with variables such as attendance rates, study hours, extracurricular involvement, and parental support. It intentionally includes data errors, missing values, and noise to simulate realistic data-processing challenges. However, exploratory analysis reveals very low correlations between most input features and the FinalGrade target variable. This lack of meaningful relationships poses a significant challenge for predictive modeling.


\section{Methodology}

\subsection{Machine Learning Model}

To predict students’ final grades, a Random Forest Regressor was selected as the primary machine learning model. Random forests are ensemble methods that construct multiple decision trees and aggregate their outputs, offering robustness against noise and overfitting. They are particularly well suited for heterogeneous tabular datasets, where relationships between variables may be nonlinear or weakly structured.

The Random Forest model was configured using the following hyperparameters:

\begin{itemize}
    \item \textbf{n\_estimators = 300}: number of trees in the ensemble, providing stable predictions.
    \item \textbf{max\_depth = None}: trees grow fully unless constrained by other stopping conditions.
    \item \textbf{min\_samples\_split = 2} and \textbf{min\_samples\_leaf = 1}: default values that allow trees to model fine-grained patterns.
    \item \textbf{max\_features = "sqrt"}: a widely used setting that increases tree diversity and reduces correlation between trees.
    \item \textbf{bootstrap = True}: enables sampling with replacement for each tree, improving generalization.
    \item \textbf{n\_jobs = -1}: uses all available processor cores to speed up training.
\end{itemize}

\subsection{Evaluation Metrics}

Model performance was assessed using three standard regression metrics that capture different aspects of prediction quality.

\subsubsection{Mean Absolute Error (MAE)}
\begin{equation}
    MAE = \frac{1}{n} \sum_{i=1}^{n} \left| y_{\text{true},i} - y_{\text{pred},i} \right|
\end{equation}
MAE measures the average magnitude of prediction errors in the original units (grade points). 
Lower values indicate better performance.

\subsubsection{Mean Squared Error (MSE) and Root Mean Squared Error (RMSE)}
\begin{equation}
    MSE = \frac{1}{n} \sum_{i=1}^{n} \left( y_{\text{true},i} - y_{\text{pred},i} \right)^{2}
\end{equation}
\begin{equation}
    RMSE = \sqrt{MSE}
\end{equation}
These metrics penalize larger errors more strongly and reflect overall prediction stability.

\subsubsection{Coefficient of Determination (R\textsuperscript{2})}
\begin{equation}
    R^{2} = 1 - \frac{\sum_{i=1}^{n} \left( y_{\text{true},i} - y_{\text{pred},i} \right)^{2}}
    {\sum_{i=1}^{n} \left( y_{\text{true},i} - \bar{y} \right)^{2}}
\end{equation}
The R\textsuperscript{2} score indicates how much variance in the target variable is explained by the model. Values close to 1 represent strong predictive power, whereas values near or below 0 indicate poor performance.

\section{Results and Discussion}

The Random Forest regression model was evaluated on the test set after completing all preprocessing steps. The results indicate very limited predictive performance. The Mean Absolute Error (MAE) is approximately {\bf7.7 grade points}, and the coefficient of determination (R²) is {\bf–0.11}, meaning the model performs worse than a simple baseline that always predicts the average final grade. This suggests that the model is unable to learn meaningful relationships between the features and the FinalGrade target.

Visualization further confirms this. The predicted values form a narrow distribution centered around the dataset mean ($\approx80$), while the actual grades span a broader range. Predicted-versus-actual plots show minimal alignment, indicating systematic underfitting. Feature importance scores are uniformly low, with no variable contributing substantially to the prediction.

These findings align with the correlation analysis, which shows near-zero correlations between all input features and the FinalGrade. The synthetic dataset—containing noise, missing values, and intentionally weak structure—does not provide the statistical signals needed for effective regression modeling. This highlights a fundamental limitation: even well-designed machine learning pipelines cannot compensate for datasets that lack meaningful predictive information.

\section{Conclusion}

This study implemented a complete data-processing and machine learning workflow to predict students’ final grades using a synthetic educational dataset. Despite thorough preprocessing and the use of a Random Forest regression model, the predictive performance remained weak. The model’s high error and negative R² score indicate that the available features do not contain sufficient information to explain or predict the FinalGrade variable.

The core limitation lies in the dataset itself. The intentionally noisy and weakly correlated synthetic features do not establish meaningful relationships with the target, making accurate regression prediction statistically infeasible. Nevertheless, the project successfully demonstrates the full pipeline of data cleaning, modeling, evaluation, and visualization.


-------------------------

Scalar {\it variables} and {\it physical constants} should be italicized, and a bold (non-italics) font should be used for {\bf vectors} and {\bf matrices}. Do not italicize subscripts unless they are variables.

Equations should be either display (with a number in parentheses) or inline. 

Display equations should be flush left and numbered consecutively, with equation numbers in parentheses and flush right.

Be sure the symbols in your equation have been defined before the equation appears or immediately following. Please refer to ``Equation (1),'' not ``Eq. (1)'' or ``equation (1).''

Punctuate display equations when they are part of the sentence preceding it, as in $a^2+b^2=c^2$. In addition, if the text following the equation flows logically as a part of the display equation, 
\begin{equation}
A=\pi r^2,
\end{equation}
use ending punctuation (comma) after the display equation.

\begin{figure}[b]
\vspace*{-10pt}
\centerline{\includegraphics[width=16pc]{fig1.png}}
\caption{Note that ``Figure'' is spelled out. There is a period after the figure number, followed by one space. It is good practice to briefly explain the significance of the figure in the caption. (Used, with permission, from \cite{AA1}.)}
\end{figure}

\begin{figure*}
\centerline{\includegraphics[width=22pc]{fig1.png}}
\caption{Note that ``Figure'' is spelled out. There is a period after the figure number, followed by one space. It is good practice to briefly explain the significance of the figure in the caption. (Used, with permission, from \cite{BB1}.)}
\end{figure*}

\section{LISTS}

Avoid using lists. Instead, use full sentences and flowing paragraphs. If you absolutely must use a list, use them rarely and keep them short:
\begin{itemize}
\item {\it Style for bulleted lists}---This is the style that should be used for bulleted lists.
	
\item {\it Punctuation in lists}---Each item in the list should end with a period, regardless of whether full sentences are used.
\end{itemize}

\section{GRAPHICAL ABSTRACTS}

This magazine accepts graphical abstracts, and they must be peer reviewed, which means the graphical abstract must be submitted with the full paper. graphical abstracts and their specifications. Please read the additional information provided by \href{http://www.ieee.org/publications_standards/publications/graphical_abstract.pdf}{{IEEE about graphical abstracts}}.\enlargethispage{7pt}


%%%%
% NOTE: The table font size must be \normalsize.
%%%%

\begin{table*}
\normalsize
\begin{center}
\caption{Units for magnetic properties.}
\label{table}
\begin{tabular*}{29.7pc}{@{}|p{29pt}|p{90pt}<{\raggedright}|p{200pt}<{\raggedright}|@{}}
\hline
Symbol& 
Quantity& 
Conversion from Gaussian and CGS EMU to SI$^{\mathrm{a}}$ \\
\hline
$\Phi $& 
Magnetic flux& 
1 Mx $\to 10^{-8}$ Wb $= 10^{-8}$ V $\cdot$ s \\
$B$& 
Magnetic flux density, magnetic induction& 
1 G $\to 10^{-4}$ T $= 10^{-4}$ Wb/m$^{2}$ \\
$H$& 
Magnetic field strength& 
1 Oe $\to 10^{-3}/(4\pi )$ A/m \\
$m$& 
Magnetic moment& 
1 erg/G $=$ 1 emu $\to 10^{-3}$ A $\cdot$ m$^{2} = 10^{-3}$ J/T \\
$M$& 
Magnetization& 
1 erg/(G $\cdot$ cm$^{3}) =$ 1 emu/cm$^{3}$ $\to 10^{-3}$ A/m \\
4$\pi M$& 
Magnetization& 
1 G $\to 10^{-3}/(4\pi )$ A/m \\
$\sigma $& 
Specific magnetization& 
1 erg/(G $\cdot$ g) $=$ 1 emu/g $\to $ 1 A $\cdot$ m$^{2}$/kg \\
$j$& 
Magnetic dipole moment& 
1 erg/G $=$ 1 emu $\to 4\pi \times 10^{-10}$ Wb $\cdot$ m \\
$J$& 
Magnetic polarization& 
1 erg/(G $\cdot$ cm$^{3}) =$ 1 emu/cm$^{3}$ $\to 4\pi \times 10^{-4}$ T \\
$\chi , \kappa $& 
Susceptibility& 
1 $\to 4\pi $ \\
$\chi_{\rho }$& 
Mass susceptibility& 
1 cm$^{3}$/g $\to 4\pi \times 10^{-3}$ m$^{3}$/kg \\
$\mu $& 
Permeability& 
1 $\to 4\pi \times 10^{-7}$ H/m $= 4\pi \times 10^{-7}$ Wb/(A $\cdot$ m) \\
$\mu_{r}$& 
Relative permeability& 
$\mu \to \mu_{r}$ \\
$w, W$& 
Energy density& 
1 erg/cm$^{3} \to 10^{-1}$ J/m$^{3}$ \\
$N, D$& 
Demagnetizing factor& 
1 $\to 1/(4\pi )$ \\
\hline
\multicolumn{3}{@{}p{29.7pc}@{}}{\vspace*{-4pt}%
\textit{Vertical lines are optional in tables. Statements that serve as captions for 
the entire table do not need footnote letters. }}\\
\multicolumn{3}{@{}p{29.7pc}@{}}{$^{\mathrm{a}}$\textit{Gaussian units are the same as cg emu for magnetostatics; Mx 
$=$ maxwell, G $=$ gauss, Oe $=$ oersted; Wb $=$ weber, V $=$ volt, s $=$ 
second, T $=$ tesla, m $=$ meter, A $=$ ampere, J $=$ joule, kg $=$ 
kilogram, H $=$ henry.}}
\end{tabular*}
\label{tab1}\vspace*{-12pt}
\end{center}
\end{table*}

\section{FIGURES AND TABLES}

\subsection{In-Text Callouts for Figures and Tables}

Figures and tables must be cited in the running text in consecutive order. At first mention, the citation should be boldface ({Figure 1}); subsequent mentions should be Roman type (see Figure 1 and {Table 1}). {Figure 2} shows an example of a figure spanning across two columns.
 
Previously published figures or tables require permission to reprint. Please obtain permission. Then, add the following text to the figure/table caption: ``From [reference no.], with permission,'' or ``Adapted from [reference no.], with permission.'' {\it Carefully} explain each figure in the text. Each manuscript should be limited to four figures.

\section{END SECTIONS}

\subsection{Appendices}

If multiple appendices are required, they should labeled ``Appendix A,'' ``Appendix B,'' etc. They appear before the ``Acknowledgment'' or the ``References'' section.

\subsection{Acknowledgment}

The ``Acknowledgment'' (no's) section appears immediately after the conclusion. If applicable, this is where you indicate funding for the work. The preferred spelling of the word ``acknowledgment'' in American English is without an ``e'' after the ``g.'' Avoid expressions such as ``One of us (S.B.A.) would like to thank $\ldots$.'' Instead, write ``We thank $\ldots$.'' Sponsor and financial support acknowledgments are included in the acknowledgment section. For example: This work was supported in part by the U.S. Department of Commerce under Grant BS123456 (sponsor and financial support acknowledgment goes here). Researchers that contributed information or assistance to the article should also be acknowledged in this section. Also, if corresponding authorship is noted in your paper it will be placed in the acknowledgment section. Note that the acknowledgment section is placed at the end of the paper before the reference section.

\subsection{References}

References need not be cited in text. When they are, they appear on the line, in square brackets, inside the punctuation. Multiple references are each\break numbered with separate brackets. When citing a section in a book, please give the relevant page numbers. In text, refer simply to the reference number. Do not use ``Ref.'' or ``reference'' except at the beginning of a sentence: ``Reference \cite{CC1} shows $\ldots$.'' Please do not use automatic endnotes in \emph{Word}, rather, type the reference list at the end of the paper using the ``References''\break style \cite{DD1}--\cite{II1}.

\enlargethispage{8pt}
Reference numbers are set flush left and form a column of their own, hanging out beyond the body of the reference. The reference numbers are on the line, end with period. In all references, the given name of the author or editor is abbreviated to the initial only and precedes the last name. Use them all; use \emph{et al.} only if names are not given. Use commas around Jr., Sr., and III in names. Abbreviate conference titles. When citing IEEE transactions, provide the issue number, page range, volume number, year, and/or month if available. When referencing a patent, provide the day and the month of issue, or application. References may not include all information; please obtain and include relevant information. Do not combine references. There must be only one reference with each number. If there is a URL included with the print reference, it can be included at the end of the reference.

Other than books, capitalize only the first word in a paper title, except for proper nouns and element symbols. For papers published in translation journals, please give the English citation first, followed by the original foreign-language citation See the end of this document for formats and examples of common\break references. For a complete discussion of references and their formats, see the IEEE style manual at \url{http://www.ieee.org/authortools/}.\vspace*{-2pt}

\section{CONCLUSION}

The manuscript should include a conclusion. In this section, summarize what was described in your paper. Future directions may also be included in this section. Authors are strongly encouraged not to reference multiple figures or tables in the conclusion; these should be referenced in the body of the paper.

\section{ACKNOWLEDGMENTS}

We thank A, B, and C. This work was supported in part by a grant from XYZ.

%%%%
% NOTE: Do not cite arXiv! Instead, give the journal titles, including vol., no., pp., and doi.
%%%%

\begin{thebibliography}{00}

\bibitem{AA1} G. M. Amdahl, G. A. Blaauw, and F. P. Brooks, ``Architecture of the IBM System/360,'' {\it IBM J. Res. \& Dev}., vol. 8, no. 2, pp. 87--101, 1964, doi: 10.1147/rd.82.0087. (journal)

\bibitem{BB1} IBM Corporation, IBM Knowledge Center - IBM Secure Service Container (Secure Service Container). [Online]. Available: {https://www.ibm.com/support/\break knowledgecenter/en/HW11R/com.ibm.hwmca.kc\_se.doc/\break introductiontotheconsole/wn2131zaci.html} (URL)

\bibitem{CC1} J. Williams, ``Narrow-band analyzer,'' PhD dissertation, Dept. of Electrical Eng., Harvard Univ., Cambridge, MA: 1993. (Thesis or dissertation)

\bibitem{DD1} J. M. P. Martinez, R. B. Llavori, M. J. A. Cabo, et al., ``Integrating data warehouses with web data: A survey,'' {\it IEEE Trans. Knowledge and Data Eng}., preprint, Dec. 21, 2007, doi: 10.1109/TKDE.2007.190746. (PrePrint)

\bibitem{EE1} W.-K. Chen, {\it Linear Networks and Systems}, Belmont, CA: Wadsworth, pp. 123--135, 1993. (book)

\bibitem{FF1} S. P. Bingulac, ``On the compatibility of adaptive controllers,'' {\it Proc. Fourth Ann. Allerton Conf. Circuits and Systems Theory}, pp. 8--16, 1994, doi: number. (conference proceedings)

\bibitem{GG1} K. Elissa, ``An overview of decision theory,'' unpublished. (Unpublished manuscript)

\bibitem{HH1} R. Nicole, ``The last word on decision theory,'' {\it J. Computer Vision}, submitted for publication. (Pending publication)

\bibitem{II1} C. J. Smith and J. S. Smith, Rocky Mountain Research Laboratories, Boulder, CO, personal communication, 1992. (Personal communication)
\end{thebibliography}%\vadjust{\vfill\pagebreak}

%%%%
% NOTE: Limit for Authors' bios is max. 4 sentences per Author.
%%%%

\newpage

\begin{IEEEbiography}{First Author}{\,}is with XYZ Corporation, New York, NY, USA. All biographies should be limited to one paragraph, four sentences, consisting of the following: (e.g., ``Dr. Author received a B.S. degree and an M.S. degree $\ldots$.''), including years achieved; association with any official journals or conferences; major professional and/or academic achievements, i.e., best paper awards, research grants, etc.; any publication information (number of papers and titles of books published); current research interests; association with any professional associations. Author membership information, e.g., is a member of the IEEE and the IEEE Consumer Technology Society, if applicable, is noted at the end of the biography. Contact him/her at faauthor@xyz.com.
\end{IEEEbiography}

\begin{IEEEbiography}{Second Author}{\,} is with ABC Corporation, B\"oblingen, Germany. Ms. Author's biography appears here. Contact him/her at sbauthor@abc.com.
\end{IEEEbiography}

\begin{IEEEbiography}{Third Author} {\,} is with DEF Corporation, Tokyo, Japan. Dr. Author's biography appears here. Contact him/her at tcauthor@def.com.
\end{IEEEbiography}

\begin{IEEEbiography}{Fourth Author} {\,} is with DEF Corporation, Tokyo, Japan. Dr. Author's biography appears here. Contact him/her at tcauthor@def.com.
\end{IEEEbiography}

\end{document}